\section{Conclusion}\label{sec:conclusion}

\para{Summary.}
We asked a simple question: what happens when you sample from a probability distribution over a hidden layer's activations the same way you sample from the final softmax, instead of passing deterministic activations forward? To answer it cleanly we built a controlled study around a single drop-in \samplinglayer, seven activation operators spanning continuous and categorical distribution types, two datasets, and 117 trained models. The answer is that it works, and in its gentle continuous form it even helps: \gaussreparam sampling costs $\approx0$ accuracy ($97.7\%$ vs $97.8\%$ on \mnist) while producing better-calibrated, lower-\nll{} models with a smaller generalization gap, and on the harder \cifar task it actually \emph{beats} the deterministic reference. The literal categorical analogue is the opposite story. Collapsing a hidden layer to a single sampled unit (\gumbelst) is a severe, high-variance information bottleneck: it is the most expensive operator on \mnist and fails to train on \cifar, collapsing to chance. This cost scales exactly as the operation becomes more faithful to a hard discrete sample --- sharpening the sample (lower $\tau$), moving the layer earlier in the network, and increasing task difficulty all make it worse. Throughout, behavior is governed first by distribution type rather than by tuning, and MC-averaging at test time recovers the single-sample cost --- precisely mirroring the well-understood behavior of sampling from a softmax. \textbf{The takeaway: sampling a hidden layer is a distribution-type-dependent regularizer, free and beneficial when continuous but a destructive bottleneck when it literally replicates one-hot softmax sampling.}

\para{Future work.}
Several directions follow directly. First, to separate the bottleneck explanation from an optimization-difficulty explanation for the \cifar collapse, temperature annealing combined with per-mode learning rates would test whether \gumbelst can be coaxed to train at all, or whether the failure is fundamentally about bits passed forward. Second, a VQ-VAE-style learned codebook \citep{oord2017vqvae} offers a less brutal categorical layer than one-hot selection, retaining discreteness while passing a richer representation. Third, multiple or structured sampling sites combined with an explicit variational information-bottleneck objective \citep{alemi2017deepvib} would let us quantify the bits each operator actually transmits, turning our empirical cost ordering into a measured information budget. Finally, scaling to deeper networks and harder data, and evaluating out-of-distribution detection and uncertainty quality from MC-sampled hidden layers, would establish whether continuous intermediate sampling is a practical route to calibrated, uncertainty-aware models rather than only a clean object of study.
